\section{Artifact and Proof-Development Notes}
\label{app:artifact}

The current manuscript is an anonymous ESOP draft in Springer LNCS format, with \texttt{main\_esop.tex} as its entry point. Before submission, the following items should be completed: (i) mechanize or fully expand the translation-correctness proof; (ii) state the preprocessing properties of FLDR as explicit lemmas and connect them to its reference implementation; (iii) decide whether the generalized $\varepsilon_r$ version belongs in the main theorem; and (iv) evaluate proof size using a stable metric, for example number of certificate cases and generated verification conditions. The anonymous author block should be retained for anonymous review.

The \texttt{synthesis/} directory contains the common kernel
(\texttt{engine.py}), problem inputs (\texttt{examples.py}), the runner
(\texttt{synthesis\_examples.py}), a pinned solver dependency, recorded
results, and explanatory notes. From that directory, run:
\begin{lstlisting}[basicstyle=\ttfamily\small]
python3 -m pip install -r requirements.txt
python3 synthesis_examples.py --output results.json
python3 diagnostics.py
python3 piecewise_examples.py
python3 piecewise_diagnostics.py
\end{lstlisting}
The default coefficient schedule is $1,2,4,8,16$, with ten seconds per solver
call and 500 candidates per input. All four recorded inputs finish with
verified certificates. \texttt{results.json} retains the coefficient vectors,
counterexamples, bound attempts, invariant checks, and separate verification
queries. \texttt{diagnostics.json} records rejection of selected invalid
certificates, rejection of a vacuous transition contract, and a direct
integer/rational audit of the concrete FLDR certificate. Exact model choices
and candidate counts can depend on solver tie-breaking.

The optional piecewise extension is implemented in the same kernel.
The runner \texttt{piecewise\_\allowbreak examples.py} reruns the four original inputs, excludes
single-affine candidates for the two-sided walk up to the configured bound,
and synthesizes certificates for that walk and three sampler inputs using
supplied partitions. \texttt{piecewise\_\allowbreak results.json} records these runs
separately from the original affine experiment. The two-sided walk's
template-wide impossibility argument is analytic, as explained in
\Cref{sec:synthesis-examples}; finite search alone does not establish it.
\texttt{piecewise\_\allowbreak diagnostics.json} additionally records rejection of
overlapping and incomplete active-state partitions and of predicates using
undeclared symbols or transition descriptors. A signed countdown checks
active-to-active region crossings in both directions. The final certificates
are reverified over the unbounded input domains, supplemented by an exact
arithmetic check on selected states.

The concrete and abstract FLDR inputs answer different questions. The former
uses the exact tree for weights $(4,5)$ and automatically includes both loop
variables $c,d$. The latter uses all declared integer symbols $m,k,n,c,d$
and a supplied structural contract for the two child outcomes. Its coverage
check ensures that admissible pairs exist; the fact that every concrete FLDR
pair is represented still follows from the stated preprocessing properties,
not from the certificate solver. The Boolean activity flag denotes the guard;
auxiliary outcome descriptors are not template features.

The artifact supports the feasibility results in
\Cref{tab:synthesis-results}; it is not a broad performance evaluation or a
mechanization of the paper's soundness and relative-completeness theorems.
The depth-only separation argument and the full-state common certificate are
both retained so that the limited scope of that separation is explicit.


For the concrete FLDR input $(4,5)$, the enabled nodes are $(0,0)$,
$(1,0)$, $(1,1)$, $(2,0)$, and $(3,0)$. The shared certificate
$W=c-2d+2$ has values $2,3,1,4,5$, respectively. The corresponding
successor-value pairs are $(3,1)$, $(4,2)$, $(0,0)$, $(5,2)$, and
$(2,0)$. Their means are $2,3,0,7/2,1$, each at most the current
value, and each pair contains a value smaller by at least one.
All enabled values are positive integers, terminal values are zero, and
$U=V=W$ establishes the domination condition. This is the direct
rational-arithmetic check recorded by the diagnostic script.
